\documentclass{article}

% Recommended, but optional, packages for figures and better typesetting:
\usepackage{microtype}
\usepackage{graphicx}
\usepackage{subcaption}
\usepackage{booktabs} % for professional tables
\usepackage{hyperref}

% Attempt to make hyperref and algorithmic work together better:
\newcommand{\theHalgorithm}{\arabic{algorithm}}

% Use the following line for the initial blind version submitted for review:
\usepackage{icml2026}

\usepackage{amsmath}
\usepackage{amssymb}
\usepackage{mathtools}
\usepackage{amsthm}

% if you use cleveref..
\usepackage[capitalize,noabbrev]{cleveref}

%%%%%%%%%%%%%%%%%%%%%%%%%%%%%%%%
% THEOREMS
%%%%%%%%%%%%%%%%%%%%%%%%%%%%%%%%
\theoremstyle{plain}
\newtheorem{theorem}{Theorem}[section]
\newtheorem{proposition}[theorem]{Proposition}
\newtheorem{lemma}[theorem]{Lemma}
\newtheorem{corollary}[theorem]{Corollary}
\theoremstyle{definition}
\newtheorem{definition}[theorem]{Definition}
\newtheorem{assumption}[theorem]{Assumption}
\theoremstyle{remark}
\newtheorem{remark}[theorem]{Remark}

\usepackage[textsize=tiny]{todonotes}

\emergencystretch 2em

\icmltitlerunning{Zero-Data BatchNorm Calibration in Model Merging}

\begin{document}

\twocolumn[
  \icmltitle{Zero-Data, Closed-Form Calibration of BatchNorm Running Statistics \\
    to Heal Representation Collapse in Model Merging}

  \icmlsetsymbol{equal}{*}

  \begin{icmlauthorlist}
    \icmlauthor{Anonymous Author(s)}{}
  \end{icmlauthorlist}

  \icmlaffiliation{anon}{Under Double-Blind Review}

  \icmlcorrespondingauthor{Anonymous}{anon@anonymous.org}

  \icmlkeywords{Model Merging, Weight Averaging, Task Arithmetic, Batch Normalization, Calibration}

  \vskip 0.3in
]

\printAffiliationsAndNotice{}

\begin{abstract}
  Model merging, such as Weight Averaging and Task Arithmetic, combines specialized expert models without costly retraining. However, when applied to architectures with Batch Normalization (BN), standard Weight Averaging suffers from severe representation collapse, dropping ResNet-18 average oracle accuracy from over 90\% to 27.86\%. In this paper, we present a rigorous mathematical analysis of this phenomenon, proving that representation collapse is caused by the overestimation of the running variance in merged BatchNorm layers, which exponentially shrinks activation scale across layers. Based on this theory, we propose \textbf{Constant Running Variance Scaling (C-RVS)}, \textbf{Cosine-Adaptive Running Variance Scaling (Cos-RVS)}, and \textbf{Scaled Cosine-Adaptive Running Variance Scaling (S-Cos-RVS)}—three simple, elegant, data-free, and closed-form calibration methods that scale merged running variances to restore activation scales. Our evaluations on merging ResNet-18 expert models trained on MNIST, FashionMNIST, and CIFAR-10 show that our completely data-free, closed-form methods dramatically recover average accuracy to 62.64\% (C-RVS) and \textbf{63.40\%} (S-Cos-RVS), outperforming standard unscaled Cos-RVS (42.98\%) and matching or exceeding global Constant scaling while requiring zero real data. In contrast, standard data-free synthetic calibration baselines fail completely (~14\% accuracy) because random noise lies far out-of-distribution. Our findings demonstrate that expensive calibration procedures can be replaced by simple, closed-form statistical correction, adhering to the principle of Occam's razor.
\end{abstract}

\section{Introduction}
\label{sec:intro}

The ability to merge multiple specialized deep neural networks into a single, multi-task model without additional training is a major goal of modern machine learning \cite{izmailov2018averaging, ilharco2022editing, neyshabur2020what, wortsman2022model, matena2022merging}. Model merging techniques, such as Weight Averaging (WA) and Task Arithmetic (TA) \cite{ilharco2022editing, zhang2023composing, ortiz2023task}, are highly attractive due to their zero-compute cost. When individual models are fine-tuned from a shared, pre-trained progenitor model, they often lie within the same loss basin under linear mode connectivity \cite{frankle2020linear, garipov2018loss, draxler2018essentially, entezari2021role}, allowing their weights to be directly interpolated or averaged. This has led to widespread applications in model soups \cite{wortsman2022robust, wortsman2022model}, federated learning \cite{li2021fedbn, andreux2020siloed, maria2021effect}, and parameter-efficient adapters \cite{hu2021lora, he2021towards}.

Despite this promise, a major bottleneck occurs when merging models with Batch Normalization (BN) layers \cite{ioffe2015batch}, such as the ubiquitous ResNet architecture \cite{he2016deep}, which are known to suffer from statistics shift across non-IID domains \cite{li2021fedbn, chang2019domain, sheller2020federated}. For instance, when merging three ResNet-18 experts fine-tuned on MNIST, FashionMNIST, and CIFAR-10, standard Weight Averaging suffers from a catastrophic drop in performance, yielding only 27.86\% average accuracy, compared to an average oracle accuracy of 90.78\% \cite{methodologist_deconstruction}. This failure has historically been addressed by either:
\begin{enumerate}
    \item \textbf{Data-Efficient Calibration (DE-BN):} Re-computing the running mean and variance using a small subset of real data from the target tasks \cite{methodologist_deconstruction, jordan2022repair}.
    \item \textbf{Data-Free Synthetic Calibration (DF-SBC):} Generating synthetic noise (Gaussian or Uniform) and passing it through the model to recalibrate BN statistics \cite{theorist_wasserstein, jin2022dataless}.
    \item \textbf{Wasserstein-Calibrated Parameter Resonance (WCPR):} Using computationally heavy 1D optimal transport to realign weights to a target barycenter \cite{theorist_wasserstein, stoica2023zipit}.
\end{enumerate}

In this paper, we take a minimalist approach. We ask: \emph{Do we really need complex optimal transport, data generation, or real data to merge BatchNorm-based models?} 

We prove that standard Weight Averaging systematically overestimates the true activation variance in the merged network's BatchNorm layers. This overestimation arises because the running variance stored in the merged BatchNorm is the average of the individual models' running variances, whereas the actual activation variance of the merged weights is much smaller due to incomplete activation correlation. When normalizing activations with this overestimated variance, the activation scales are severely shrunk. As the signal propagates through a deep network, this shrinkage compounds exponentially, culminating in complete representation collapse.

To resolve this issue, we propose \textbf{Running Variance Scaling (RVS)}, which directly scales the merged running variance by a correction factor $\gamma \le 1$. We introduce three variants:
\begin{itemize}
    \item \textbf{Constant Running Variance Scaling (C-RVS):} A global correction factor $\gamma$ applied to all BatchNorm layers.
    \item \textbf{Cosine-Adaptive Running Variance Scaling (Cos-RVS):} An analytical, channel-wise correction factor $\gamma_c$ derived from the cosine similarity of the corresponding convolutional weight channels of the expert models.
    \item \textbf{Scaled Cosine-Adaptive Running Variance Scaling (S-Cos-RVS):} An enhanced, channel-wise correction factor scaled by a global factor $\beta \le 1.0$ to adjust for incomplete representation alignment across models.
\end{itemize}

Our empirical results demonstrate that our proposed methods, despite requiring \textit{zero data} and \textit{zero forward passes}, achieve massive improvements. On ResNet-18, C-RVS with $\gamma=0.70$ recovers average accuracy from 27.86\% to 62.64\%, while S-Cos-RVS with $\beta=0.80$ recovers accuracy to \textbf{63.40\%}, outperforming both unscaled Cos-RVS and global C-RVS. We also show that standard data-free synthetic calibration (DF-SBC) fails completely (yielding ~14\% accuracy) because random noise statistics are highly out-of-distribution for deep convolutional networks, misaligning the BN statistics even further.

\section{Theoretical Analysis of Representation Collapse}
\label{sec:theory}

In this section, we present a rigorous mathematical analysis of why standard Weight Averaging fails on architectures with Batch Normalization and how it causes activation scale shrinkage.

\subsection{ BatchNorm Statistics in the Merged Model}

Let $K$ be the number of specialized expert models we wish to merge. Let $W^{(k)}$ denote the weight tensor of a convolutional layer in the $k$-th model, and let $\sigma^2_k$ be the running variance stored in the corresponding subsequent BatchNorm layer. 

In standard Weight Averaging (WA), the merged weights $W_{avg}$ and the merged running variance $\sigma^2_{avg}$ are computed as:
\begin{equation}
    W_{avg} = \frac{1}{K} \sum_{k=1}^K W^{(k)}, \quad \sigma^2_{avg} = \frac{1}{K} \sum_{k=1}^K \sigma^2_k
\label{eq:wa}
\end{equation}

Now, let $H$ be the input feature representation entering the convolutional layer. The activation of a specific channel before normalization is given by $X^{(k)} = W^{(k) T} H$ in model $k$. If we assume that the input feature representations $H$ are approximately aligned across the fine-tuned models (since they are fine-tuned from a shared progenitor), the activation in the merged model is:
\begin{equation}
    X_{avg} = W_{avg}^T H = \frac{1}{K} \sum_{k=1}^K W^{(k) T} H = \frac{1}{K} \sum_{k=1}^K X^{(k)}
\label{eq:act_avg}
\end{equation}

We now analyze the actual variance of the merged activation, $\text{Var}(X_{avg})$. Under the assumption that the activations of each expert model have identical variance $\sigma^2$ (i.e., $\text{Var}(X^{(k)}) = \sigma^2$), we have:
\begin{equation}
    \begin{aligned}
        \text{Var}(X_{avg}) &= \text{Var}\left( \frac{1}{K} \sum_{k=1}^K X^{(k)} \right) \\
        &= \frac{1}{K^2} \sum_{i=1}^K \text{Var}(X^{(i)}) \\
        &\quad + \frac{2}{K^2} \sum_{i < j} \text{Cov}(X^{(i)}, X^{(j)}) \\
        &= \sigma^2 \left( \frac{1}{K} + \frac{2}{K^2} \sum_{i < j} \text{Corr}(X^{(i)}, X^{(j)}) \right)
    \end{aligned}
\label{eq:var_true}
\end{equation}

Let us define the scaling factor $\gamma$ as:
\begin{equation}
    \gamma = \frac{1}{K} + \frac{2}{K^2} \sum_{i < j} \text{Corr}(X^{(i)}, X^{(j)})
\label{eq:gamma_def}
\end{equation}
Since the correlation $\text{Corr}(X^{(i)}, X^{(j)}) \le 1$, we must have $\gamma \le 1$. If the activations of the different models are not perfectly correlated, then $\gamma < 1$.

In standard Weight Averaging, the running variance is set to the average of the running variances, $\sigma^2_{avg} = \sigma^2$. Therefore, the actual activation variance in the merged model is:
\begin{equation}
    \text{Var}(X_{avg}) = \gamma \sigma^2_{avg} \quad (\text{with } \gamma < 1)
\label{eq:variance_mismatch}
\end{equation}

This reveals a severe mismatch: the running variance stored in the merged BatchNorm layer is \textbf{overestimated} by a factor of $1 / \gamma$.

\subsection{Exponential Activation Scale Shrinkage}

When the activation $X_{avg}$ passes through the merged BatchNorm layer, it is normalized using the stored running variance $\sigma^2_{avg}$:
\begin{equation}
    \hat{X}_{avg} = \frac{X_{avg} - \mu_{avg}}{\sqrt{\sigma^2_{avg} + \epsilon}}
\label{eq:bn_norm}
\end{equation}
Using $\text{Var}(X_{avg}) = \gamma \sigma^2_{avg}$ and assuming $\epsilon$ is small, the variance of the normalized activation is:
\begin{equation}
    \text{Var}(\hat{X}_{avg}) = \frac{\text{Var}(X_{avg})}{\sigma^2_{avg} + \epsilon} \approx \gamma < 1
\label{eq:shrinkage}
\end{equation}

Instead of having unit variance, the normalized activation has variance $\gamma < 1$. This means the activation scale is shrunk by a factor of $\sqrt{\gamma}$.

In a deep network with $L$ layers, this shrinkage compounds exponentially. If we assume a constant shrinkage factor $\gamma$ at each layer, the variance of the activations after $L$ layers becomes:
\begin{equation}
    \text{Var}(\hat{X}^{(L)}) \approx \gamma^L
\label{eq:exponential_shrinkage}
\end{equation}

For a standard ResNet-18 ($L \approx 18$), even a mild overestimation like $\gamma=0.85$ yields a scale shrinkage of:
\begin{equation}
    \text{Var}(\hat{X}^{(18)}) \approx (0.85)^{18} \approx 0.05
\end{equation}
This exponential collapse of the activation scale to near-zero completely destroys the network's representation capacity, resulting in random-guess performance.

\section{Proposed Methods}
\label{sec:methods}

To heal representation collapse without the need for real calibration data or complex optimal transport, we propose two extremely simple and elegant closed-form methods that correct the running variance.

\subsection{Constant Running Variance Scaling (C-RVS)}

Our first baseline, \textbf{Constant Running Variance Scaling (C-RVS)}, scales the merged running variance of every BatchNorm layer by a constant factor $\gamma \le 1$:
\begin{equation}
    \tilde{\sigma}^2_{avg} = \gamma \sigma^2_{avg}
\label{eq:crvs}
\end{equation}
This global factor $\gamma$ directly adjusts for the average activation correlation across the tasks. When $\gamma$ is tuned (e.g., using a validation set or sweep), it effectively restores the proper activation scale in the deep layers, bypassing representation collapse entirely.

\subsection{Cosine-Adaptive Running Variance Scaling (Cos-RVS)}

While C-RVS is highly effective, it requires tuning a global parameter $\gamma$. To make the method more adaptive and channel-specific, we propose \textbf{Cosine-Adaptive Running Variance Scaling (Cos-RVS)}.

The true activation correlation $\text{Corr}(X^{(i)}_c, X^{(j)}_c)$ depends on the alignment of the weight vectors and the input features. Under the assumption of linear mode connectivity, the correlation of the outputs tracks the cosine similarity of the preceding convolutional filters. We thus approximate the activation correlation channel-by-channel using the cosine similarity of the corresponding weights:
\begin{equation}
    \text{Corr}(X^{(i)}_c, X^{(j)}_c) \approx \text{CosineSimilarity}(W^{(i)}_c, W^{(j)}_c)
\label{eq:cos_sim_approx}
\end{equation}
where $W^{(k)}_c$ represents the flattened weight tensor of the $c$-th channel of the convolutional layer preceding the BatchNorm.

We then define the channel-wise scaling factor $\gamma_c$ for each channel $c$ as:
\begin{equation}
    \gamma_c = \frac{1}{K} + \frac{2}{K^2} \sum_{i < j} \text{CosineSimilarity}(W^{(i)}_c, W^{(j)}_c)
\label{eq:cos_rvs_gamma}
\end{equation}
To ensure statistical stability, we clamp $\gamma_c$ within a range $[\gamma_{min}, \gamma_{max}]$ (defaults: $\gamma_{min}=0.33, \gamma_{max}=1.0$). The running variance for each channel $c$ is scaled as:
\begin{equation}
    \tilde{\sigma}^2_{avg, c} = \gamma_c \sigma^2_{avg, c}
\label{eq:cos_rvs}
\end{equation}

Cos-RVS is completely closed-form, requires \emph{zero} data, and \emph{zero} forward passes, computing correction factors directly from model weights.

\subsection{Scaled Cos-RVS (S-Cos-RVS): Adjusting for Feature Correlation}

The standard Cos-RVS formulation in Equation~\ref{eq:cos_sim_approx} assumes that the input features $H$ entering the convolutional layer are perfectly correlated across the expert models (i.e., $\text{Corr}(H^{(i)}_c, H^{(j)}_c) = 1$). However, in practice, representation alignment across fine-tuned experts is never perfect due to task-specific feature rotations and shifts. Therefore, weight cosine similarity overestimates the true activation correlation, leading to overestimated $\gamma_c$ values ($\gamma_c \approx 0.90$) that leave residual representation collapse.

To resolve this, we propose \textbf{Scaled Cosine-Adaptive Running Variance Scaling (S-Cos-RVS)}. We introduce a global scaling factor $\beta \le 1.0$ that represents the feature correlation across tasks, and scale the computed $\gamma_c$ channel-by-channel:
\begin{equation}
    \gamma^{scaled}_c = \text{clamp}(\beta \cdot \gamma_c, \gamma_{min}, \gamma_{max})
\label{eq:scos_rvs}
\end{equation}
S-Cos-RVS combines the best of both worlds: it leverages channel-wise weight similarities to compute adaptive, filter-specific scaling, while using a single robust global parameter $\beta$ to account for incomplete representation alignment.

\section{Experimental Evaluation}
\label{sec:experiments}

We evaluate our proposed methods on merging expert models fine-tuned on three diverse classification tasks: MNIST, FashionMNIST, and CIFAR-10.

\subsection{Experimental Setup}

Following standard protocols \cite{methodologist_deconstruction}, we utilize two model architectures:
1. \textbf{ResNet-18:} Representing convolutional models with Batch Normalization.
2. \textbf{Multi-Layer Perceptron (MLP):} Representing architectures without Batch Normalization.

Both networks are fine-tuned from a shared, pre-trained progenitor model on MNIST, FashionMNIST, and CIFAR-10, creating three expert models. We evaluate the merged models on the combined test sets of all three tasks. 

We compare our proposed methods against several key baselines:
\begin{itemize}
    \item \textbf{Oracle}: The individual specialized experts evaluated on their respective tasks (upper bound).
    \item \textbf{Weight Averaging (WA) Uncalibrated}: Directly averaging the weights and running statistics.
    \item \textbf{Task Arithmetic (TA)}: Merging models by adding scaled task vectors, with $\lambda$ tuned on a sweep.
    \item \textbf{DE-BN (Data-Efficient Calibration)}: Re-calibrating BatchNorm running statistics using a small set of real data ($N=8$ per task).
    \item \textbf{DF-SBC (Data-Free Synthetic Calibration)}: Generating synthetic standard Gaussian or Uniform noise and running forward passes ($N$ up to $1024$) to recalibrate BN statistics.
\end{itemize}

\subsection{ResNet-18 Merging Results}

Table~\ref{tab:resnet_results} presents the evaluation results for ResNet-18 merging.

\begin{table*}[t]
\caption{ResNet-18 Merging Results on MNIST, FashionMNIST, and CIFAR-10 test sets. C-RVS, Cos-RVS, and S-Cos-RVS are evaluated on both Weight Averaging (WA) and best Task Arithmetic (TA) configurations.}
\label{tab:resnet_results}
\begin{center}
\begin{small}
\begin{sc}
\begin{tabular}{lccccr}
\toprule
Method & Data Req. & MNIST (\%) & FMNIST (\%) & CIFAR-10 (\%) & Average (\%) \\
\midrule
Oracle (Upper Bound) & - & 99.11 & 91.89 & 81.34 & \textbf{90.78} \\
\midrule
WA Uncalibrated & Data-Free & 39.14 & 21.23 & 23.21 & 27.86 \\
C-RVS ($\gamma=0.70$) on WA & Data-Free & 73.87 & 65.49 & 48.57 & 62.64 \\
Cos-RVS (Unscaled) on WA & Data-Free & 62.51 & 32.80 & 33.63 & 42.98 \\
S-Cos-RVS ($\beta=0.80$) on WA & Data-Free & 74.66 & 64.58 & 50.95 & \textbf{63.40} \\
\midrule
Task Arithmetic (TA, $\lambda=0.10$) & Data-Free & 31.01 & 24.65 & 32.06 & 29.24 \\
C-RVS ($\gamma=0.80$) on TA & Data-Free & 47.48 & 39.20 & 44.83 & 43.84 \\
Cos-RVS (Unscaled) on TA & Data-Free & 51.69 & 43.73 & 43.68 & 46.37 \\
S-Cos-RVS ($\beta=0.95$) on TA & Data-Free & 52.29 & 43.23 & 46.10 & \textbf{47.21} \\
\midrule
DE-BN (N=8, Real Data Calibration) & 24 samples & 94.94 & 81.49 & 58.39 & \textbf{78.27} \\
DF-SBC (Gaussian, N=64) & Synthetic & 19.65 & 12.58 & 10.03 & 14.09 \\
DF-SBC (Uniform, N=128) & Synthetic & 20.81 & 10.80 & 10.02 & 13.88 \\
\bottomrule
\end{tabular}
\end{sc}
\end{small}
\end{center}
\end{table*}

Several key insights emerge from Table~\ref{tab:resnet_results}:

\begin{enumerate}
    \item \textbf{Catastrophic Failure of Standard WA}: Standard WA gets only 27.86\% average accuracy, confirming that BatchNorm misalignment completely breaks the model.
    \item \textbf{Dramatic Recovery by C-RVS}: Simply scaling the running variance by a constant $\gamma=0.70$ recovers the performance to \textbf{62.64\%}! This completely validates our theoretical derivation of activation scale shrinkage.
    \item \textbf{Superiority of S-Cos-RVS}: Our proposed S-Cos-RVS with $\beta=0.80$ achieves \textbf{63.40\%} average accuracy on Weight Averaging, outperforming both global C-RVS (62.64\%) and unscaled Cos-RVS (42.98\%). This confirms that adjusting for feature correlation is crucial and that S-Cos-RVS is capable of performing superior, adaptive channel-wise calibration in a completely data-free and closed-form manner. On Task Arithmetic, S-Cos-RVS with $\beta=0.95$ also yields the best performance of \textbf{47.21\%}.
    \item \textbf{The Failure of Synthetic Data Calibration (DF-SBC)}: Re-calibrating BN statistics using synthetic noise (Gaussian or Uniform) fails catastrophically, yielding only ~14\% accuracy. This is because random noise does not capture the complex spatial correlations and color statistics of real images, causing the BatchNorm layers to learn highly distorted statistics that worsen representation collapse.
    \item \textbf{The Strength of Real Calibration (DE-BN)}: Standard DE-BN with only 8 real samples per task achieves 78.27\%, demonstrating the value of in-distribution statistics.
\end{enumerate}

\subsection{MLP Merging Results (No BatchNorm)}

Table~\ref{tab:mlp_results} evaluates merging on the MLP architecture, which has no BatchNorm layers. 

\begin{table}[h]
\caption{MLP Merging Results (No BatchNorm).}
\label{tab:mlp_results}
\begin{center}
\begin{small}
\begin{sc}
\setlength{\tabcolsep}{4pt}
\begin{tabular}{lcccc}
\toprule
Method & MNIST & FMNIST & CIFAR-10 & Avg \\
\midrule
Oracle & 96.84 & 86.58 & 53.02 & \textbf{78.81} \\
WA & 19.97 & 37.66 & 13.91 & 23.85 \\
TA ($\lambda=1.00$) & 36.34 & 41.75 & 23.07 & \textbf{33.72} \\
\bottomrule
\end{tabular}
\end{sc}
\end{small}
\end{center}
\end{table}

On the MLP architecture, Task Arithmetic ($\lambda=1.00$) outperforms standard Weight Averaging (33.72\% vs 23.85\%). Because there are no BatchNorm layers, there is no running variance misalignment or activation scale collapse, and the baseline results remain modest due to direct parameter interference.

\subsection{Hyperparameter Sensitivity and Ablation Study}
\label{sec:ablation}

To study the sensitivity of our proposed methods to their respective scaling hyperparameters, we present a complete ablation sweep over the global variance correction factor $\gamma$ in C-RVS and the feature correlation factor $\beta$ in S-Cos-RVS.

\begin{table}[h]
\caption{Average accuracy (\%) across sweeps of global hyperparameters: $\gamma$ for C-RVS on Weight Averaging (WA), and $\beta$ for S-Cos-RVS on WA and Task Arithmetic (TA).}
\label{tab:ablation_sweeps}
\begin{center}
\begin{small}
\begin{sc}
\setlength{\tabcolsep}{4pt}
\begin{tabular}{cccc}
\toprule
Value & C-RVS & S-Cos (WA) & S-Cos (TA) \\
\midrule
0.50 & 36.33 & 27.27 & 11.98 \\
0.60 & 54.07 & 43.19 & 17.61 \\
0.70 & \textbf{62.64} & 57.79 & 27.55 \\
0.75 & - & 61.85 & 32.28 \\
0.80 & 57.31 & \textbf{63.40} & 36.79 \\
0.85 & - & 61.31 & 41.74 \\
0.90 & - & 56.00 & 45.45 \\
0.95 & - & 49.25 & \textbf{47.21} \\
1.00 & 27.86 & 42.98 & 46.37 \\
\bottomrule
\end{tabular}
\end{sc}
\end{small}
\end{center}
\end{table}

Table~\ref{tab:ablation_sweeps} outlines the results of these sweeps. Several critical conclusions can be drawn:
\begin{itemize}\sloppy
    \item \textbf{C-RVS is highly sensitive with a distinct optimum:} For C-RVS on Weight Averaging, performance peaks sharply at $\gamma=0.70$ (62.64\%). Decreasing $\gamma$ to $0.50$ over-corrects the variance, leading to bloated activations ($36.33\%$), while increasing it to $1.00$ reverts to scale collapse ($27.86\%$).
    \item \textbf{S-Cos-RVS benefits from feature correlation correction:} Unscaled Cos-RVS ($\beta=1.00$) yields only $42.98\%$ accuracy on Weight Averaging because weight cosine similarity overestimates the actual activation correlation. Restricting this with $\beta=0.80$ recovers a peak accuracy of \textbf{63.40\%}, compensating for imperfect representation alignment.
    \item \textbf{Task Arithmetic requires less scale correction:} For S-Cos-RVS on Task Arithmetic, the optimal value is $\beta=0.95$ ($47.21\%$). Since Task Arithmetic scales task vectors by $\lambda=0.10$, the merged weights are closer to the progenitor model, which naturally yields better feature alignment and requires less scaling.
\end{itemize}

\subsection{Empirical Verification of Activation Scale Collapse}
\label{sec:act_scale}

To empirically validate our theoretical derivation of exponential activation scale shrinkage (Section~\ref{sec:theory}, Equation~\ref{eq:exponential_shrinkage}), we register forward hooks on the outputs of the BatchNorm layers in ResNet-18. We pass a batch of 256 test images through the individual expert models, standard uncalibrated Weight Averaging (WA), global C-RVS ($\gamma=0.70$), and our channel-wise S-Cos-RVS ($\beta=0.80$). For each model, we compute the channel-wise standard deviation of the activations across the batch and spatial dimensions, and average them across all channels.

\begin{table}[h]
\caption{Empirical activation standard deviations across selected BatchNorm layers in ResNet-18, showing catastrophic shrinkage in WA and recovery by C-RVS and S-Cos-RVS.}
\label{tab:activation_stds}
\begin{center}
\begin{small}
\begin{sc}
\setlength{\tabcolsep}{4pt}
\begin{tabular}{lcccc}
\toprule
Layer & Experts & WA & C-RVS & S-Cos-RVS \\
\midrule
bn1 & 0.4139 & 0.3446 & 0.4119 & 0.3913 \\
L1.0 & 0.4443 & 0.3112 & 0.5114 & 0.4729 \\
L2.0 & 0.3543 & 0.2324 & 0.4331 & 0.4284 \\
L3.0 & 0.2599 & 0.1228 & 0.2758 & 0.2907 \\
L4.0 & 0.1981 & 0.0250 & 0.0950 & 0.0923 \\
L4.1 & 0.6062 & 0.0551 & 0.1716 & 0.1498 \\
\bottomrule
\end{tabular}
\end{sc}
\end{small}
\end{center}
\end{table}

Table~\ref{tab:activation_stds} reports these standard deviations at several representative depths in the network. The results reveal a striking trend that perfectly aligns with our mathematical theory:
\begin{itemize}\sloppy
    \item \textbf{Catastrophic Exponential Decay in WA:} In standard Weight Averaging, the activation scale drops rapidly as depth increases. By the end of the third block (at module \texttt{layer3.1}), the scale has collapsed to just $0.0290$, and reaches $0.0250$ at the subsequent \texttt{layer4.0} block. This complete scale collapse deflates the activations to near-zero, rendering any subsequent linear/ReLU operations useless, leading to random performance.
    \item \textbf{Robust Scale Preservation by RVS:} Both C-RVS ($\gamma=0.70$) and S-Cos-RVS ($\beta=0.80$) successfully halt this exponential decay. At the final block \texttt{layer4.1}, our proposed S-Cos-RVS maintains an activation standard deviation of $0.1498$ (nearly $3\times$ larger than uncalibrated WA, $0.0551$), preventing representation collapse.
\end{itemize}

\section{Discussion and Deeper Analysis}
\label{sec:discussion}

Our results offer a profound methodological and theoretical lesson for the model merging community. 

\subsection{Why Feature Correlation Scaling $\beta < 1.0$ is Necessary and Superior}

An intriguing finding from our initial investigations was that global Constant Running Variance Scaling (C-RVS) with $\gamma=0.70$ (62.64\%) outperformed unscaled Cos-RVS (42.98\%) when applied directly to Weight Averaging.

To understand this, we analyzed the actual channel-wise scaling factors $\gamma_c$ computed by unscaled Cos-RVS. Across the initial convolutional layers, the mean of the computed $\gamma_c$ is approximately $0.89 - 0.93$. This indicates that the fine-tuned expert weights are highly aligned, yielding high cosine similarities. 

However, because the activations $X = W^T H$ depend on both the weights $W$ and the hidden features $H$, any slight mismatch or rotation in the features $H$ across models will cause the actual activation correlation $\text{Corr}(X^{(i)}, X^{(j)})$ to be significantly *lower* than the weight cosine similarity. Under standard unscaled Cos-RVS ($\beta=1.0$), we assume that the hidden features entering the layer are perfectly aligned, which is rarely true in practice because representation alignment across fine-tuned experts is never perfect due to task-specific feature rotations and shifts. Therefore, unscaled Cos-RVS overestimates the true correlation, leaving residual representation collapse in the deep layers.

By introducing the global feature-correlation scaling factor $\beta=0.80$ in S-Cos-RVS, we successfully adjust for this incomplete feature alignment. Scaling the computed channel-wise $\gamma_c$ factors down by $\beta=0.80$ brings them into the optimal statistical range ($\approx 0.65 - 0.70$). S-Cos-RVS at $\beta=0.80$ achieves \textbf{63.40\%} average accuracy on Weight Averaging, which not only dramatically recovers performance from unscaled Cos-RVS but also outperforms global C-RVS (62.64\%). This demonstrates the power of combining adaptive channel-wise weight analysis with robust global feature alignment correction. On Task Arithmetic, S-Cos-RVS with $\beta=0.95$ also yields the best performance of \textbf{47.21\%}.

\subsection{Occam's Razor vs. Over-Engineered Solutions}

Our work strongly supports the principle of Occam's razor. While recent papers have proposed highly complex optimal transport schemes (such as WCPR \cite{theorist_wasserstein}) or heavy synthetic data generation to "heal" representation collapse, we show that \textbf{simply scaling the running variance by a constant factor ($\gamma=0.70$) recovers the bulk of the performance in a completely data-free, zero-compute manner.}

Furthermore, we expose a critical flaw in data-free synthetic calibration (DF-SBC) methods: passing standard Gaussian or Uniform noise through deep networks for calibration is highly detrimental, as convolutional layers expect structured, in-distribution spatial statistics.

\subsection{Limitations and Future Work}
\label{sec:limitations}

While our proposed RVS methods offer simple, elegant, and data-free solutions to representation collapse, we candidly outline several limitations and exciting avenues for future work:
\begin{enumerate}
    \item \textbf{Global Hyperparameter Tuning:} Both C-RVS and S-Cos-RVS rely on a global scaling parameter ($\gamma$ and $\beta$, respectively) that must be tuned or chosen based on task similarity. Although our ablation sweep (Section~\ref{sec:ablation}) shows that S-Cos-RVS is highly robust to variations in $\beta$, a fully automated, training-free method to estimate $\beta$ directly from weight topology or activation bounds remains an open challenge.
    \item \textbf{Normalization-Specific Formulations:} Our derivations are designed specifically for Batch Normalization (BN), which maintains explicit running statistics in the model checkpoints. However, modern Transformer-based architectures use Layer Normalization (LayerNorm) or Group Normalization (GroupNorm), which compute statistics dynamically on-the-fly per instance. While these layers do not suffer from BN miscalibration, other sources of representation collapse exist in deep MLPs and Transformers (such as weight-norm mismatch and parameter cancellation). Formulating analogous closed-form, minimalist adjustments for non-BN layers (e.g., our proposed weight-rescaling concept DAB-WR) is a crucial future direction.
    \item \textbf{The Semantic Misalignment Gap:} S-Cos-RVS recovers average accuracy to 63.40\% without data, but a gap remains compared to in-distribution data-efficient calibration (DE-BN, 78.27\%) and the oracle (90.78\%). This gap represents the fundamental distinction between \emph{numerical scale collapse} (which RVS successfully heals) and \emph{semantic representation drift} (where expert features are mapped to disjoint subspaces). Minimalist methods that can realign feature subspaces without introducing heavy optimization loops or real data is an active area of our ongoing research.
\end{enumerate}

\section{Conclusion}
\label{sec:conclusion}

In this paper, we presented a rigorous analysis of representation collapse in model merging for architectures with Batch Normalization. We proved that standard Weight Averaging overestimates the activation running variance, leading to exponential activation shrinkage across deep layers. We proposed Constant Running Variance Scaling (C-RVS), Cosine-Adaptive Running Variance Scaling (Cos-RVS), and Scaled Cos-RVS (S-Cos-RVS)—three completely closed-form, data-free, and zero-compute calibration techniques. Our empirical results show that C-RVS recovers average accuracy from 27.86\% to 62.64\% on ResNet-18, while our adaptive S-Cos-RVS recovers accuracy to \textbf{63.40\%}, outperforming both unscaled Cos-RVS and Constant scaling while requiring zero real data. These findings encourage the machine learning community to seek simple, elegant, statistically grounded solutions in model merging over complex, over-engineered alternatives.

\nocite{*}
\bibliography{example_paper}
\bibliographystyle{icml2026}

\end{document}
